\documentclass[11pt]{article}
\usepackage[utf8]{inputenc}
\usepackage[T1]{fontenc}
\usepackage{amsmath,amssymb}
\usepackage{booktabs}
\usepackage{graphicx}
\usepackage[margin=1in]{geometry}
\usepackage[numbers]{natbib}
\usepackage{hyperref}

\title{When Does Recursion Pay Off?\\
\large An Open, \$0 Reproduction and Honest Dissection of\\
Two-Timescale Meta-Skill Evolution}
\author{loversky02 \\ \texttt{github.com/loversky02/metaskill-evolve}}
\date{2026}

\newcommand{\metaskill}{\textsc{MetaSkill-Evolve}}

\begin{document}
\maketitle

\begin{abstract}
Self-improving LLM agents rewrite their own skill files from execution traces.
\metaskill{}~\citep{wang2026metaskill} makes this \emph{recursive} with a
two-timescale scheme: a fast loop evolves a task skill $s$, while a slow loop
evolves a branch-local meta-skill $m=(\psi,\sigma,\alpha,\pi,\varepsilon)$ that
parameterises the five agents of the improvement pipeline itself, under the same
pipeline applied to itself. We contribute (i) an open, \$0, Apple-Silicon
reproduction of the method with a single frozen backbone, and (ii) an honest
dissection of \emph{when} the recursion pays off. Decomposing the original
results shows the fast loop is the workhorse and the recursive meta-loop
contributes a modest, task-dependent increment ($+6.4/+8.1/+1.9$ points over
single-level evolution on OfficeQA/SealQA/ALFWorld). We hypothesise a
\emph{backbone-strength law}: the meta-loop's marginal value shrinks as the frozen
backbone strengthens, making recursive self-improvement chiefly worthwhile for
cheap or weak models. We instrument this with a controlled mechanism study and
two further probes---English$\rightarrow$Vietnamese meta transfer, and the token
cost of recursion. We are explicit about scope: absolute numbers require the
paper's 31B backbone, which does not fit our 24\,GB budget; our engine reproduces
the \emph{method} and the \emph{qualitative structure}, and live confirmation of
the backbone-strength law on real backbones is ongoing.
\end{abstract}

\section{Introduction}
External skills---reusable procedural knowledge handed to an agent---extend what
LLM agents can do on long-horizon tasks. Self-improving agents rewrite these skill
files from their own traces. \metaskill{} observes that such self-evolution is
still \emph{non-recursive}: it improves \emph{what} the agent does (the task
skill) but not \emph{how} it improves (the improvement procedure), which is
authored once and frozen. \metaskill{} closes this gap with a second, slower loop
that evolves the improvement procedure with the same machinery.

This paper is not a new method. It is a reproduction and a dissection, asking a
single question a practitioner cares about: \textbf{given a fixed backbone and a
token budget, is the recursion worth it, and when?} Our contributions:
\begin{itemize}
  \item An open, dependency-light engine (\S\ref{sec:engine}) that reproduces the
  two-timescale method and runs entirely offline at \$0 on a 24\,GB Mac.
  \item A decomposition showing the fast loop dominates and the recursion is a
  modest, task-dependent bonus (\S\ref{sec:workhorse}).
  \item The backbone-strength law as a testable hypothesis, with a controlled
  mechanism study (\S\ref{sec:law}), plus EN$\rightarrow$VI transfer
  (\S\ref{sec:transfer}) and a cost analysis (\S\ref{sec:cost}).
\end{itemize}

\section{Background: \metaskill{}}
Let $U(s)=\mathbb{E}_{(x,y)\sim\mathcal{T}}[r(A_s(x),y)]$ be the utility of a task
skill $s$, where $A_s$ is the agent conditioned on $s$ and $r\in[0,1]$ scores
predictions. The fast loop (Alg.~1) mines a worst case, runs the pipeline
Analyzer $\to$ Retriever $\to$ Allocator $\to$ Proposer $\to$ Evolver to produce
$K$ children, scores them, and commits to a branch DAG. Every $H$ iterations the
slow loop (Alg.~2) aggregates meta-productivity
$P(m\mid s)=\mathbb{E}\!\left[\tfrac{1}{K}\sum_{k}\big(U(s'_k)-U(s)\big)\right]$,
synthesises a meta-failure trace, and applies the \emph{same} pipeline to the five
meta files $\{\psi,\sigma,\alpha,\pi,\varepsilon\}$. All five agents share one
frozen backbone. Key defaults: $H{=}2$, $K_{\max}{=}3$, frontier size $3$ with
weights $(1,0.5,0.25)$.

\section{An Open, \$0 Reproduction}\label{sec:engine}
Our engine mirrors the method: versioned Markdown skill/meta-skill files with
snapshot/restore, a SQLite branch DAG with a weighted frontier, the five agents
reading their meta files, and a provider-agnostic backbone (mock for tests, a
local MLX model for the weak end, or an OpenAI-compatible gateway for the strong
end). The whole pipeline is text-only with a frozen backbone---no gradient
updates---so both loops run at \$0. The reference implementation is 33 unit tests
green with no API key, GPU, or downloads.

\section{Dissection}
\subsection{The fast loop is the workhorse}\label{sec:workhorse}
Decomposing Table~1 of \citet{wang2026metaskill} into the single-level jump (fast
loop, meta frozen) and the recursive increment (Table~\ref{tab:decomp}) shows that
$70$--$90\%$ of the gain over the raw backbone comes from the fast loop alone; the
recursion adds $+6.4/+8.1/+1.9$. On ALFWorld the baseline is already $92.3\%$
(near ceiling; static skills even \emph{hurt}), so the $+1.9$ is within noise and
a poor test of recursion.

\begin{table}[h]\centering
\caption{Decomposition of the reported gains. Recursion is the small part.}
\label{tab:decomp}
\begin{tabular}{lccccc}
\toprule
Benchmark & No-skill & Static & Single-level & \metaskill{} & Recursion $\Delta$\\
\midrule
OfficeQA & 31.78 & 36.09 & 48.94 & 55.32 & $+6.38$\\
SealQA   & 29.17 & 29.41 & 37.21 & 45.26 & $+8.05$\\
ALFWorld & 92.31 & 90.38 & 92.31 & 94.23 & $+1.92$\\
\bottomrule
\end{tabular}
\end{table}

\subsection{A backbone-strength law}\label{sec:law}
\textbf{Hypothesis.} The recursive meta-loop rewrites the prompts that scaffold
the pipeline agents. A strong backbone already analyses and proposes well, so this
scaffolding is redundant; a weak backbone benefits more. We therefore predict the
meta-gain \emph{decreases} as the backbone strengthens---so recursive
self-improvement is chiefly worthwhile for cheap/weak models, which is also the
economically interesting regime.

\textbf{Mechanism study (toy).} We construct a controlled task where an agent's
behaviour reads a directive from its own meta file, so the slow loop can causally
improve the pipeline. With proxy backbones of increasing competence, the meta-gain
falls monotonically to zero (Table~\ref{tab:law-toy}). \emph{These are toy numbers
that verify the mechanism and the harness, not benchmark results.}

\textbf{Live probe: both failure modes bracket a narrow regime.} We ran the
procedure on real backbones---Claude Haiku~4.5 and Opus~4.8 (gateway) and a local
Gemma-2-2B (MLX). On short multi-step word problems, every backbone reached
\emph{ceiling} before any evolution ($U_0{=}1.0$): nothing to improve, meta-gain
$0$. We then tried a genuinely hard benchmark, SealQA Seal-Hard~\citep{pham2025sealqa},
with a bare LLM (no tools); there Haiku scored $0/12$ ($U_0{=}0.0$)---the opposite
\emph{floor} failure. The cause is instructive: SealQA is a \emph{knowledge/search}
benchmark (fresh, trap-laden fact-seeking), and evolving a text skill on a frozen
backbone cannot supply missing knowledge---the paper's own SealQA gains came from an
agent \emph{with} search tools, which our \$0 bare-LLM reproduction deliberately
omits. Between these two failure modes lies the narrow regime where recursive
skill-evolution actually pays off: tasks whose bottleneck is an \emph{improvable
procedure}, not triviality (ceiling) or missing knowledge/tools (floor).
\emph{Characterising that regime---rather than assuming recursion always
helps---is the honest empirical contribution here.} Table~\ref{tab:law-live}
records all four live failure modes.

\textbf{Two escape attempts, two more modes.} To move off ceiling/floor we built
two extensions (both released): a keyless search-augmented agent (so SealQA becomes
a reason-over-noisy-results task, not a knowledge test), and a synthetic
\emph{procedure-bottleneck} task (letter counting, where a systematic-enumeration
skill should help). Probing real backbones exposed two further modes.
\emph{(iii) Too weak to self-improve}: a local Gemma-2-2B has genuine headroom on
letter counting ($U_0{=}0.33$), yet the fast and slow loops leave it unchanged
($+0.00$)---the same 2B that must \emph{author} the skill cannot write a useful one.
\emph{(iv) Noise-dominated}: Claude Haiku on a harder count has headroom
($U_0{\approx}0.4$--$0.6$) and is capable, but on a small stochastic eval the loop
chases sampling noise and can \emph{degrade} the skill (a run gave single/two-level
$0.4/0.0$ against $U_0{=}0.6$). A clean positive live signal thus needs the paper's
fuller regime \emph{simultaneously}: a capable backbone, tools/context, and a large
low-variance eval---which our deliberately minimal \$0 bare-LLM reproduction does
not assemble. The offline mechanism study stays the controlled demonstration; live,
we offer a taxonomy of failure modes, which we think is the more honest deliverable.

\begin{table}[h]\centering
\caption{Mechanism study on a toy with proxy backbones (offline). Meta-gain
vanishes as the backbone strengthens.}
\label{tab:law-toy}
\begin{tabular}{lccc}
\toprule
Backbone (proxy) & Single-level & Two-level & Meta-gain\\
\midrule
weak   & 0.50 & 1.00 & $+0.50$\\
mid    & 0.75 & 1.00 & $+0.25$\\
strong & 1.00 & 1.00 & $+0.00$\\
\bottomrule
\end{tabular}
\end{table}

\begin{table}[h]\centering
\caption{Live probe: four failure modes bracket the narrow regime where recursive
skill-evolution helps. A clean positive needs headroom AND a capable-enough backbone
AND a low-variance eval \emph{at once} --- which a minimal \$0 bare-LLM setup does
not assemble.}
\label{tab:law-live}
\begin{tabular}{llccl}
\toprule
Task & Backbone & $U_0$ & Two-level & Failure mode\\
\midrule
Word problems     & haiku / opus / Gemma-2-2B  & 1.00 & 1.00 & ceiling (trivial)\\
SealQA Seal-Hard  & Claude Haiku (no tools)    & 0.00 & ---  & floor (no knowledge)\\
Letter-count      & Gemma-2-2B (MLX)           & 0.33 & 0.33 & too weak to author skills\\
Letter-count (hard) & Claude Haiku 4.5         & 0.60 & 0.00 & noise-dominated (small eval)\\
\bottomrule
\end{tabular}
\end{table}

\subsection{English$\rightarrow$Vietnamese meta transfer}\label{sec:transfer}
Meta-skills are language-agnostic improvement policies. We ask whether a meta-skill
evolved on English tasks accelerates task-skill evolution on Vietnamese ones (GSM8K
$\rightarrow$ vi-gsm8k). The harness runs; live numbers are \textsc{TODO}.

\subsection{The cost of recursion}\label{sec:cost}
Recursion multiplies inference---you evolve the evolvers. We instrument calls and
tokens and compare fast-only vs.\ two-level at equal iteration budget, reporting
utility gained per token. On the toy the recursion costs roughly $+30\%$ tokens
for its extra gain; the live token/accuracy trade-off is \textsc{TODO}.

\section{Related Work}
Skill-augmented and self-evolving agents are a fast-growing cluster
\citep{skillrl2026,evolvingrl2026}. \metaskill{} is distinctive in making the
improvement procedure itself recursive with no extra model. Our work is orthogonal:
we do not propose a new evolver but dissect \emph{when} recursion helps, on a
frozen backbone, at \$0.

\section{Limitations and Honesty}
(1) The paper's absolute numbers used a 31B backbone (~20\,GB at 4-bit) that does
not fit our 24\,GB budget; we reproduce the method and qualitative structure, not
the exact numbers. (2) The backbone-strength law is currently supported only by a
constructed mechanism study; the proxy backbones were designed to express the
predicted shape, so they demonstrate the harness, not the law. (3) OfficeQA data is
access-gated; SealQA and ALFWorld are public. Live experiments are the immediate
next step.

\section{Conclusion}
Two-timescale recursion is an elegant idea, but its benefit is modest and
task-dependent, concentrated where the improvement pipeline---not the base
policy---is the bottleneck. We predict, and are testing, that this benefit is a
decreasing function of backbone strength. An open \$0 engine makes the question
cheap for anyone to probe.

\bibliographystyle{plainnat}
\bibliography{refs}
\end{document}
